\documentclass[12pt]{article}
\usepackage[utf8]{inputenc}
\usepackage[T1]{fontenc}
\usepackage[margin=1in]{geometry}
\usepackage{setspace}
\usepackage{xcolor}
\usepackage[colorlinks=true, urlcolor=blue!60!black, linkcolor=blue!60!black]{hyperref}
\setlength{\parindent}{2em}
\setlength{\parskip}{1.25em}
\renewcommand{\baselinestretch}{1.0}
\begin{document}
\begin{center}
{\LARGE \textit{Baker v. General Motors Corp.}, 522 U.S. 222 (1998)\par}
\vspace{0.8em}
{\large Case Brief --- Four-Page Edition\par}
\end{center}
\vspace{1em}

\section{Memory Jogger}

A Michigan consent injunction barring a former GM engineer from testifying against GM could not stop him from testifying, under subpoena, in a Missouri wrongful-death suit brought by strangers to the Michigan case: sister-state judgments get exacting full faith and credit with no public-policy exception, but enforcement measures do not travel and a decree cannot control other courts' litigation (\href{https://supreme.justia.com/cases/federal/us/522/222/}{522 U.S. 222}).

\section{Detailed Case Facts}

For fifteen years Ronald Elwell worked as a GM engineering analyst studying vehicle fires and often assisted GM's lawyers in product liability defense. After the relationship soured, he sued GM in Michigan for wrongful discharge and GM counterclaimed. In 1992 they settled: GM paid an undisclosed sum, and the Michigan court entered a stipulated permanent injunction barring Elwell from testifying against GM without its consent in any litigation involving GM's products. The decree carved out a pending Georgia case, and a side agreement said that testimony ordered by another tribunal would ``in no way'' support a contempt or breach-of-contract action.

Beverly Garner died in a February 1990 Missouri crash and fire while riding in a 1985 Chevrolet S-10 Blazer.

\section{Procedural History}

Her sons, Kenneth Lee and Steven Robert Baker, sued GM in Missouri state court alleging a faulty fuel pump; GM removed to federal court. The Bakers, strangers to the Michigan case, subpoenaed Elwell. The district court allowed his testimony, reasoning that enforcing the gag order would offend Missouri's policy favoring disclosure and that the injunction was modifiable. The jury awarded \$11.3 million. The Eighth Circuit reversed, and the Supreme Court granted certiorari.

\section{Judicial Votes}

Reversed, \textbf{9--0} in the judgment, with a \textbf{5--4} split on rationale.
\begin{itemize}
\item \textbf{Majority:} Justice Ginsburg, joined by Chief Justice Rehnquist and Justices Stevens, Souter, and Breyer.
\item \textbf{Concurring in the judgment:} Justice Scalia; Justice Kennedy, joined by Justices O'Connor and Thomas.
\item \textbf{Dissents:} None.
\end{itemize}

\section{Holding}

Elwell could testify in the Missouri action without offending the Full Faith and Credit Clause or 28 U.S.C. \S~1738. The full faith and credit owed final judgments is exacting and admits no roving public-policy exception, unlike the choice-of-law treatment of statutes. But a sister-state equitable decree cannot dictate evidentiary rules or override another jurisdiction's subpoena power in litigation by strangers over whom the rendering court lacked authority.

\section{Analysis of Opinions}

\textbf{Ginsburg (majority).} \textit{Laws versus judgments.} The Clause makes the States ``integral parts of a single nation.'' A forum need not apply a sister state's statutes, but a judgment rendered with jurisdiction commands nationwide preclusion, even on a claim the forum would not enforce, as in \textit{Fauntleroy v.\ Lum}. \textit{Equity and enforcement.} Equitable decrees are not categorically exempt, yet the forum need not adopt the rendering state's ``time, manner, and mechanisms for enforcing judgments.'' \textit{Limits on injunctions.} Decrees lack extraterritorial effect when they purport to perform an official act within another state's province, like transferring title to land (\textit{Fall v.\ Eastin}), or interfere with litigation beyond the rendering court's authority. Michigan validly barred Elwell from volunteering as an expert, but it had no jurisdiction over the Bakers (\textit{Martin v.\ Wilks}). The Georgia carve-out and GM's side agreement showed GM was using Missouri as a ``surrogate tribunal.''

\textbf{Scalia (concurring in the judgment).} Under \textit{McElmoyle} and \textit{Thompson v.\ Whitman}, full faith and credit is a rule of evidence and preclusion, not execution: a judgment ``cannot be enforced out of the state by an execution issued within it.''

\textbf{Kennedy (concurring in the judgment).} The majority's exceptions sit in ``uneasy tension'' with its rejection of a public-policy exception. Under \textit{Kremer} and \textit{Matsushita}, the question is the preclusive effect in the rendering state, and Michigan requires mutuality and does not bind nonparties.

\textbf{Later developments.} \textit{Hyatt I} (2003) relied on the laws--judgments divide; \textit{Lance v.\ Dennis} (2006) cited \textit{Baker} on nonparty preclusion; \textit{V.L.\ v.\ E.L.} (2016) invoked its exacting rule to require recognition of a Georgia adoption decree.

\section{Examples: Future Applications}

\begin{enumerate}
\item \textbf{Same-Side: A Gagged Scientist.} State~A bars a former pharmaceutical scientist from testifying about side effects. Nonparty plaintiffs in State~B may subpoena her, because State~A cannot control State~B's subpoenas.
\item \textbf{Same-Side: Garnishment Rules.} A State~A money judgment is registered in State~B, which caps garnishment and exempts homesteads. State~B must recognize the debt, but its own enforcement rules govern.
\item \textbf{Opposite-Side: Money Judgment on a Settlement.} A State~A money judgment for breaching a non-disclosure settlement must be enforced in State~B despite contrary policy, because it only fixes a debt between parties.
\item \textbf{Opposite-Side: The Bakers Lose in Michigan.} Had the Bakers intervened and lost a privilege ruling in Michigan, Missouri would owe it issue-preclusive effect.
\item \textbf{Fence-Sitter: Shield Laws.} State~A enters a money judgment against a State~B physician for prescribing abortion medication by telehealth, and State~B's shield law forbids enforcement. \textit{Baker}'s exacting rule favors enforcement, but penal-judgment and personal-jurisdiction defenses remain unsettled.
\end{enumerate}

\section{Critique}

\textbf{Against the majority.} Kennedy's concurrence exposes the core vulnerability in Ginsburg's opinion: if judgments admit no roving public-policy exception, carving out injunctions that interfere with sister-state litigation risks smuggling a functional public-policy exception in through a jurisdictional label.

\textbf{Against Kennedy.} Ginsburg's Footnote~12 answers that GM claimed only to enforce an injunction against Elwell, a Michigan party, not to preclude the Bakers, so state mutuality doctrine alone cannot explain why Missouri may compel Elwell to defy a valid Michigan order.

\section{Key Quotations}

\begin{quote}
``Regarding judgments, however, the full faith and credit obligation is exacting. A final judgment in one State, if rendered by a court with adjudicatory authority over the subject matter and persons governed by the judgment, qualifies for recognition throughout the land. \ldots But our decisions support no roving `public policy exception' to the full faith and credit due judgments.''
\end{quote}
\noindent --- Justice Ginsburg, Opinion of the Court, 522 U.S. at 233.

\begin{quote}
``Full faith and credit, however, does not mean that States must adopt the practices of other States regarding the time, manner, and mechanisms for enforcing judgments. Enforcement measures do not travel with the sister state judgment as preclusive effects do; such measures remain subject to the evenhanded control of forum law.''
\end{quote}
\noindent --- Justice Ginsburg, Opinion of the Court, 522 U.S. at 235.

\begin{quote}
``Most essentially, Michigan lacks authority to control courts elsewhere by precluding them, in actions brought by strangers to the Michigan litigation, from determining for themselves what witnesses are competent to testify and what evidence is relevant and admissible in their search for the truth.''
\end{quote}
\noindent --- Justice Ginsburg, Opinion of the Court, 522 U.S. at 238.

\end{document}
